\chapter{Appendix}

\section{Appendix A: Important Source Code Excerpts}
This appendix lists selected implementation excerpts from important VisuaLearn AI features. The complete project contains additional route handlers, services, components, and configuration files; only representative code is included here to show the core implementation approach.

\subsection{Central API Route Registration}
The backend keeps feature routes modular and registers them from a central API module. This helps separate chat, learning content, documents, feedback, simulation, visualization, and Notion integration.

\begin{tcblisting}{listing only,colback=gray!10!white,breakable,boxsep=0pt,top=1mm,bottom=1mm,left=1mm,right=1mm,
listing options={basicstyle=\footnotesize\ttfamily,backgroundcolor=\color{gray!20},breaklines=true,columns=fullflexible}}
// backend/src/api/index.js
import express from 'express';
import learningContentRoutes from './learning/index.js';
import documentRoutes from './documents/routes.js';
import feedbackRoutes from './feedback/routes.js';
import simulationRoutes from './simulation/routes.js';
import visual3dRoutes from './visual3d/routes.js';
import notionRoutes from './notion/routes.js';
import chatRoutes from '../../routes/chat.js';

export function registerRoutes(app) {
  const router = express.Router();

  router.use('/learning-content', learningContentRoutes);
  router.use('/documents', documentRoutes);
  router.use('/feedback', feedbackRoutes);
  router.use('/simulation', simulationRoutes);
  router.use('/visual3d', visual3dRoutes);
  router.use('/notion', notionRoutes);
  router.use('/chat', chatRoutes);

  app.use('/api', router);
}
\end{tcblisting}

\subsection{Document Upload and RAG Preparation}
Uploaded PDFs are validated, stored, recorded in the database, and then processed into chunks for document-based learning.

\begin{tcblisting}{listing only,colback=gray!10!white,breakable,boxsep=0pt,top=1mm,bottom=1mm,left=1mm,right=1mm,
listing options={basicstyle=\footnotesize\ttfamily,backgroundcolor=\color{gray!20},breaklines=true,columns=fullflexible}}
// backend/src/api/documents/routes.js
router.post('/upload', upload.single('file'), async (req, res) => {
  const userId = req.user?.userId;
  if (!userId) return res.status(401).json({ error: 'Authentication required' });
  if (!req.file) return res.status(400).json({ error: 'No file provided' });

  const validation = validatePdf(req.file.buffer);
  if (!validation.valid) {
    return res.status(400).json({ error: validation.error });
  }

  const filename = `${uuidv4()}.pdf`;
  const storagePath = `${userId}/${filename}`;

  await supabase.storage.from('documents').upload(storagePath, req.file.buffer, {
    contentType: 'application/pdf',
    upsert: false,
  });

  const document = await createDocument(
    userId,
    filename,
    req.file.originalname,
    req.file.size,
    storagePath,
    req.file.mimetype
  );

  await processPdfBuffer(document.id, req.file.buffer);
  res.json({ success: true, documentId: document.id });
});
\end{tcblisting}

\subsection{Embedding Endpoint Builder}
The RAG service supports both a full Azure embeddings URL and a base endpoint. This makes the embedding configuration flexible while still routing requests through the backend.

\begin{tcblisting}{listing only,colback=gray!10!white,breakable,boxsep=0pt,top=1mm,bottom=1mm,left=1mm,right=1mm,
listing options={basicstyle=\footnotesize\ttfamily,backgroundcolor=\color{gray!20},breaklines=true,columns=fullflexible}}
// backend/src/services/rag/index.js
export function buildEmbeddingRequestUrl(endpoint, deployment, apiVersion) {
  if (!endpoint) return null;

  const parsed = new URL(endpoint.trim());

  if (
    parsed.pathname.includes('/openai/deployments/') &&
    parsed.pathname.endsWith('/embeddings')
  ) {
    if (!parsed.searchParams.has('api-version')) {
      parsed.searchParams.set('api-version', apiVersion);
    }
    return parsed.toString();
  }

  const basePath = parsed.pathname.replace(/\/+$/, '');
  parsed.pathname = `${basePath}/openai/deployments/${deployment}/embeddings`;
  parsed.search = '';
  parsed.searchParams.set('api-version', apiVersion);
  return parsed.toString();
}
\end{tcblisting}

\subsection{Simulation API}
Simulation requests are routed through detection and generation endpoints. The response returns a structured simulation payload instead of exposing arbitrary frontend code execution.

\begin{tcblisting}{listing only,colback=gray!10!white,breakable,boxsep=0pt,top=1mm,bottom=1mm,left=1mm,right=1mm,
listing options={basicstyle=\footnotesize\ttfamily,backgroundcolor=\color{gray!20},breaklines=true,columns=fullflexible}}
// backend/src/api/simulation/routes.js
router.post('/detect', rateLimitSimulation, async (req, res) => {
  const { query } = req.body;
  if (!query || typeof query !== 'string') {
    return res.status(400).json({ success: false, error: 'query is required' });
  }

  const decision = LearningOrchestratorDecision({
    query,
    mode: req.body.mode || 'chat',
    conversationState: req.body.conversationState || {},
    requestedArtifact: req.body.requestedArtifact || null,
  });

  const support = detectSandboxSimulationSupport(decision.activeTopic || query, {
    explicitQuery: query,
    requestedArtifact: req.body.requestedArtifact || null,
  });

  res.json({ success: true, supported: support.supported, decision, ...support });
});
\end{tcblisting}

\subsection{Lazy Learning Content Loading}
The frontend learning hook fetches only the requested learning tab. This prevents every artifact from being generated for every normal conversation message.

\begin{tcblisting}{listing only,colback=gray!10!white,breakable,boxsep=0pt,top=1mm,bottom=1mm,left=1mm,right=1mm,
listing options={basicstyle=\footnotesize\ttfamily,backgroundcolor=\color{gray!20},breaklines=true,columns=fullflexible}}
// frontend/src/hooks/useLearningContent.js
const fetchTabContent = useCallback(async (
  query,
  tabType,
  userId = null,
  accessToken = null,
  preferences = null,
  conversationId = null,
  documentId = null
) => {
  if (!query?.trim() || !tabType) return null;

  const headers = { 'Content-Type': 'application/json' };
  if (accessToken) headers.Authorization = `Bearer ${accessToken}`;

  const response = await fetch(`${API_BASE}/api/learning-content`, {
    method: 'POST',
    headers,
    body: JSON.stringify({
      query,
      userId,
      contentType: tabType,
      preferences,
      conversationId,
      documentId,
    }),
  });

  return parseJsonResponse(response, 'Learning content service unavailable.');
}, []);
\end{tcblisting}

\section{Glossary}
\begin{description}[leftmargin=3.2cm,style=nextline]
	\item[Learning Artifact] A generated resource such as notes, quiz, flashcards, mind map, simulation, or visualization.
	\item[Adaptive Strategy] The selected teaching format for a learner response.
	\item[RAG] Retrieval-augmented generation using uploaded document chunks.
	\item[Simulation] A step-by-step representation of a procedural topic.
	\item[Persistence] Storage of conversations and generated resources for future access.
\end{description}

\clearpage
\section{Appendix B: Selected System Screenshots}
This appendix includes selected screenshots from the implemented VisuaLearn AI application. The figures show the main user-facing workflows across chat, persona selection, learning content, quizzes, flashcards, mind maps, simulation, and 3D visualization.

\begin{figure}[H]
\centering
\includegraphics[width=0.96\textwidth]{Figures/normal.png}
\caption{Chat mode interface showing an active conversation with streamed assistant response, message history, and input controls.}
\label{fig:appendix-normal-chat-entry}
\end{figure}

\begin{figure}[H]
\centering
\includegraphics[width=0.96\textwidth]{Figures/chat.png}
\caption{New conversation screen showing the primary VisuaLearn AI entry point for starting a general learning or chat session.}
\label{fig:appendix-chat-mode}
\end{figure}

\begin{figure}[H]
\centering
\includegraphics[width=0.96\textwidth]{Figures/persona.png}
\caption{Persona selection interface used to choose the tutoring style that guides the assistant's tone and explanation approach.}
\label{fig:appendix-persona-selection}
\end{figure}

\begin{figure}[H]
\centering
\includegraphics[width=0.96\textwidth]{Figures/learn.png}
\caption{Learning mode screen presenting structured explanatory content for a selected topic inside the learning workspace.}
\label{fig:appendix-learn-mode}
\end{figure}

\begin{figure}[H]
\centering
\includegraphics[width=0.96\textwidth]{Figures/quiz.png}
\caption{Quiz view generated for a learning topic, allowing the learner to test understanding through interactive questions.}
\label{fig:appendix-quiz-view}
\end{figure}

\begin{figure}[H]
\centering
\includegraphics[width=0.96\textwidth]{Figures/flashcards.png}
\caption{Flashcards view showing concise question-and-answer cards for revision and repeated practice.}
\label{fig:appendix-flashcards-view}
\end{figure}

\begin{figure}[H]
\centering
\includegraphics[width=0.96\textwidth]{Figures/mindmaps.png}
\caption{Mind map view organizing the selected topic into connected concepts for visual exploration.}
\label{fig:appendix-mindmap-view}
\end{figure}

\begin{figure}[H]
\centering
\includegraphics[width=0.96\textwidth]{Figures/sim.png}
\caption{Simulation view showing an interactive step-by-step representation of a learning concept.}
\label{fig:appendix-simulation-view}
\end{figure}

\begin{figure}[H]
\centering
\includegraphics[width=0.96\textwidth]{Figures/3d.png}
\caption{3D visualization view demonstrating an interactive spatial representation generated for a learning concept.}
\label{fig:appendix-3d-visualization}
\end{figure}
